\section{摘要}
\label{sec:abstract}

多功能雷达/电子对抗网络在现代电子战中同时承担探测感知与电子干扰双重任务，其空间部署直接影响战场电磁态势。然而，复杂区域雷达部署面临两个核心挑战：一是区域边缘探测概率显著低于中心区域（文献报道下降15\%--20\%）的边界效应问题；二是空中节点与地面节点的传播特性差异显著（A2G链路路径损耗指数约2.0，G2G链路约4.0），但现有方法通常采用统一路径损耗模型，与实际传播环境不符。

针对上述问题，本文基于MOPSO-DT算法框架，提出面向空地协同电子对抗网络的多目标部署优化方法。首先采用Hertel-Mehlhorn算法对复杂部署区域进行凸分解，利用垂直交点法将归一化坐标映射到合法物理部署区域，以抑制边界效应；其次引入空地异构传播模型，区分空-地链路（路径损耗指数$\alpha=2.0$）和地-地链路（$\alpha=4.0$）的传播特性；最后通过多目标粒子群优化生成覆盖率与干扰压制效能的Pareto最优权衡方案。在参数调优中，采用standard惯性权重策略和crowding拥挤度全局最优选择，显著提升了Pareto前沿质量。

在300km$\times$300km区域、8部雷达的雷达方程模型场景下，算法生成6个Pareto最优解，有效覆盖率（ECR）范围为81.0\%--89.0\%，优化耗时211.9秒。在200km$\times$200km、衰减系数$\beta=0.03$的挑战性场景下，生成15个Pareto最优解，ECR范围为3.6\%--20.4\%。两类场景的结果显示，ECR与最小等效干扰强度$J_{\text{min}}$之间呈强负相关（$r<-0.93$），说明在有限节点数量和距离衰减约束下，覆盖范围扩展与最薄弱任务点压制强度之间存在明显权衡。

\textbf{关键词：}多目标粒子群优化；区域分解；坐标变换；空地协同；电子对抗；Pareto最优

\vspace{12pt}
\noindent\rule{\linewidth}{0.5pt}
\vspace{12pt}

\section*{Abstract}

We study the deployment optimization of multi-function radar/electronic-countermeasure networks in complex regions. Two challenges motivate this work: the boundary effect where detection probability degrades near region edges (reported 15--20\% drop in prior work), and the heterogeneous propagation difference between air-to-ground (A2G, $\alpha \approx 2.0$) and ground-to-ground (G2G, $\alpha \approx 4.0$) links that existing methods model uniformly.

We extend the MOPSO-DT framework with three key components: (1) Hertel-Mehlhorn convex decomposition with vertical-intersection coordinate transformation to suppress boundary effects; (2) a heterogeneous air-ground propagation model distinguishing A2G links ($\alpha=2.0$) from G2G links ($\alpha=4.0$); and (3) multi-objective particle swarm optimization with standard inertia and crowding-based global-best selection for improved Pareto diversity.

In a 300km$\times$300km scenario with 8 radars under a radar-equation model, our method produces 6 Pareto-optimal solutions with effective coverage rate (ECR) ranging 81.0\%--89.0\% in 211.9 seconds. In a challenging 200km$\times$200km scenario ($\beta=0.03$), we obtain 15 Pareto solutions with ECR ranging 3.6\%--20.4\%. The two tested scenarios show a strong negative correlation between ECR and the minimum equivalent jamming strength $J_{\text{min}}$ ($r < -0.93$), indicating an empirical trade-off between coverage expansion and weakest-point suppression strength under finite-node and distance-attenuation constraints.

\textbf{Keywords:} multi-objective particle swarm optimization; region decomposition; coordinate transformation; air-ground coordination; electronic countermeasure; Pareto optimality
